\documentclass[11pt]{book}

%%%%%%%%%%%%%%%%%%%%%%%%%%%%%%%%%%%%%%%%%%%%%%%%%%%%%%%%%%%%%%%%%%%%%%%%%%%%%%%%%%%%%%%%%%%%%%%%%%%
%                                                                                                 %
% The mathematical style of these documents follows                                               %
%                                                                                                 %
% A. Thompson and B.N. Taylor. The NIST Guide for the Use of the International System of Units.   %
%    NIST Special Publication 881, 2008.                                                          %
%                                                                                                 %
% http://www.nist.gov/pml/pubs/sp811/index.cfm                                                    %
%                                                                                                 %
%%%%%%%%%%%%%%%%%%%%%%%%%%%%%%%%%%%%%%%%%%%%%%%%%%%%%%%%%%%%%%%%%%%%%%%%%%%%%%%%%%%%%%%%%%%%%%%%%%%

\input{../../fds/Manuals/Bibliography/commoncommands}
\IfFileExists{gitrevision.tex}
{\input{gitrevision}}
{\newcommand{\gitrevision}{unknown} }

% Math shortcuts
%\renewcommand{_{\rm }[1]{_\mathrm{#1}}

\addbibresource{../../fds/Manuals/Bibliography/FDS_general.bib}
\addbibresource{../../fds/Manuals/Bibliography/FDS_mathcomp.bib}
\addbibresource{../../fds/Manuals/Bibliography/FDS_refs.bib}

\begin{document}

\pagestyle{empty}

\begin{minipage}[t][9in][s]{6.25in}

\headerA{
1169\\
}

\headerB{
Verification and Validation\\
of Commonly Used\\
Empirical Correlations\\
for Fire Scenarios\\
}

\normalsize

\headerC{
\flushright{Kristopher J. Overholt}
}

\vfill

\headerD{1169}

\vfill

\end{minipage}

\newpage
\hspace{5in}
\newpage

\frontmatter

\pagenumbering{roman}

\begin{minipage}[t][9in][s]{6.25in}

\headerA{
1169\\}

\headerB{
Verification and Validation\\
of Commonly Used\\
Empirical Correlations\\
for Fire Scenarios\\
}

\headerC{
\flushright{
Kristopher J. Overholt \\
{\em Fire Research Division \\
Engineering Laboratory} } }


\headerD{1169}

\headerC{
\flushright{\today \\
Revision:~\gitrevision}}

\vfill

\flushright{\includegraphics[width=1in]{../../fds/Manuals/Bibliography/doc} }

\titlesigs

\end{minipage}

\newpage


\disclaimer{1169}


\clearpage


\pagestyle{plain}
\pagenumbering{roman}

\chapter{Preface}

In 2007, the U.S. Nuclear Regulatory Commission (NRC), together with the Electric Power Research Institute (EPRI) and the National Institute of Standards and Technology (NIST), conducted a research project to verify and validate five fire models that have been used for nuclear power plant (NPP) applications. The results of this effort were documented in a seven-volume report, NUREG-1824 (EPRI 1011999), Verification and Validation of Selected Fire Models for Nuclear Power Plant Applications~\cite{NUREG_1824}.

In 2014, the verification and validation study was expanded, and this document was created to serve as a verification and validation guide for the empirical correlations. The full details of this expanded verification and validation study are summarized in NUREG-1824 Supplement 1 (EPRI 3002002182)~\cite{NUREG_1824_Sup_1}.

The model evaluation process consists of two main components: verification and validation. Verification is a process to check the correctness of the solution of the governing equations. Verification does not imply that the governing equations are appropriate; only that the equations are being solved correctly.

Validation is a process to determine the appropriateness of the governing equations as a mathematical model of the physical phenomena of interest. Typically, validation involves comparing model results with experimental measurement. Differences that cannot be attributed to uncertainty in the measured quantities in the experiment are attributed to the assumptions and simplifications of the physical model.

\input{../../fds/Manuals/Bibliography/disclaimer}

\cleardoublepage
\phantomsection
\addcontentsline{toc}{chapter}{Contents}
\tableofcontents

\cleardoublepage
\phantomsection
\addcontentsline{toc}{chapter}{List of Figures}
\listoffigures

\cleardoublepage
\phantomsection
\addcontentsline{toc}{chapter}{List of Tables}
\listoftables

\chapter{List of Acronyms}

\begin{tabbing}
\hspace{1.5in} \= \\
AST \> Adiabatic Surface Temperature \\
ATF \> Bureau of Alcohol, Tobacco, Firearms, and Explosives \\
CAROLFIRE \> Cable Response to Live Fire Test Program \\
CFAST \> Consolidated Model of Fire Growth and Smoke Transport \\
FDS \> Fire Dynamics Simulator \\
FM \> Factory Mutual Global \\
HGL \> Hot Gas Layer \\
HRR \> Heat Release Rate \\
LLNL \> Lawrence Livermore National Laboratory \\
NBS \> National Bureau of Standards (former name of NIST) \\
NFPRF \> National Fire Protection Research Foundation \\
NIST \> National Institute of Standards and Technology \\
NRC \> Nuclear Regulatory Commission \\
RTI \> Response Time Index \\
SNL \> Sandia National Laboratory \\
SP \>  Statens Provningsanstalt (Technical Research Institute of Sweden) \\
THIEF \> Thermally-Induced Electrical Failure \\
UL  \> Underwriters Laboratories \\
USN \> United States Navy \\
VTT \> Valtion Teknillinen Tutkimuskeskus (Technical Research Centre of Finland) \\
WTC \> World Trade Center \\
\end{tabbing}

\mainmatter

\include{Introduction_Chapter}

\include{HGL_Temperature_Chapter}

\include{HGL_Depth_Chapter}

\include{Plume_Temperature_Chapter}

\include{Target_Temperature_Chapter}

\include{Target_Heat_Flux_Chapter}

\include{Ceiling_Jet_Temperature_Chapter}

\include{Sprinkler_Activation_Time_Chapter}

\include{Smoke_Detector_Activation_Time_Chapter}

\include{Summary_Chapter}

\nopart %To Fix TOC in PDF output.
\printbibliography[heading=bibintoc,title={References}]

\appendix

\include{Inputs_Chapter}

\end{document}
